
\documentclass[11pt]{article}
\usepackage[ngerman]{babel}
\usepackage[a4paper,margin=0.9in]{geometry}
\usepackage[T1]{fontenc}
\usepackage[utf8]{inputenc}
\usepackage{lmodern}
\usepackage{amsmath,amssymb,mathtools,mathrsfs}
\usepackage{microtype}
\usepackage{fancyhdr}
\usepackage{array,longtable,enumitem,multicol}
\usepackage{tikz}
\setlength{\parindent}{1.5em}
\setlength{\parskip}{0.18em}
\emergencystretch=6em
\pagestyle{fancy}
\fancyhf{}
\cfoot{\thepage}
\newcommand{\Jac}[2]{\left(\frac{#1}{#2}\right)}
\newcommand{\disc}{\operatorname{disc}}
\newcommand{\Norm}{\operatorname{N}}
\newcommand{\Tr}{\operatorname{Tr}}
\providecommand{\dd}{\,d}
\providecommand{\eps}{\varepsilon}
\providecommand{\ii}{\mathrm{i}}
\providecommand{\Om}{\Omega}
\providecommand{\frakm}{\mathfrak m}
\providecommand{\fraka}{\mathfrak a}
\providecommand{\frakb}{\mathfrak b}
\providecommand{\frakp}{\mathfrak p}
\providecommand{\calO}{\mathcal O}
\lhead{Weber Vol. I \S 39}\rhead{German source}
\begin{document}

\section*{\S\,39. Algorithmus zur Berechnung der Wurzeln.}

Der im vorigen Paragraphen gegebene Beweis für die Existenz einer Wurzel einer algebraischen Gleichung lässt zwar an Bündigkeit nichts zu wünschen übrig, er hat aber noch den Mangel, dass er nicht die Schritte erkennen lässt, wie man eine Wurzel durch ein convergentes Rechnungsverfahren berechnen kann. Wir fügen also noch die folgenden Betrachtungen hinzu, die diesem Mangel, wenn auch nur theoretisch, abzuhelfen bestimmt sind. Sie geben noch einen zweiten Beweis des Fundamentalsatzes, der in der Hauptsache von Lipschitz herrührt (Lehrbuch der Analysis, Bd. I, \S\,61 ff.; eine Vereinfachung verdanke ich einer Mittheilung von Dedekind).

Wenn die beiden ganzen rationalen Functionen $f(x)$ und $f'(x)$ einen gemeinsamen Theiler haben, so kann dieser nach \S\,6 durch rationale Operationen gefunden und beseitigt werden. Wir dürfen also voraussetzen, dass $f(x)$ und $f'(x)$ keinen gemeinsamen Theiler haben, dass sie also für keinen Werth von $x$ beide zugleich verschwinden.

Wir setzen nun voraus, dass die Wurzelberechnung für eine Function $(n-1)$ten Grades schon gelungen sei, und nehmen demnach $f'(x)$ in lineare Factoren zerlegt an. Demnach sei, wenn
\begin{equation}
 f(x)=x^n+a_1x^{n-1}+a_2x^{n-2}+\cdots,
\tag{1}
\end{equation}
ist,
\begin{equation}
 f'(x)=nx^{n-1}+(n-1)a_1x^{n-2}+\cdots
      =n(x-\beta_1)(x-\beta_2)\cdots(x-\beta_{n-1}),
\tag{2}
\end{equation}
worin die $\beta_1,\beta_2,
\ldots,\beta_{n-1}$ als bekannte Zahlen angesehen werden, die auch theilweise identisch sein können. Nach unserer Voraussetzung werden die absoluten Werthe
\begin{equation}
 |f(\beta_1)|,\ |f(\beta_2)|,\ldots, |f(\beta_{n-1})|,
\tag{3}
\end{equation}
die wir mit
\begin{equation}
 b_1,b_2,\ldots,b_{n-1}
\tag{4}
\end{equation}
bezeichnen, alle von Null verschieden sein; wir wollen die Bezeichnung so gewählt annehmen, dass $b_1$ der kleinste unter ihnen sei, oder wenigstens keinen der anderen an Grösse übertrifft. Nach dem Satz 5., \S\,38 lässt sich dann ein Werth $\alpha$ von $x$ so bestimmen, dass der absolute Werth $a$ von $f(\alpha)$ kleiner als $b_1$ wird, also:
\begin{equation}
 a<b_1,b_2,\ldots,b_{n-1}.
\tag{5}
\end{equation}

Wir begrenzen nun ein Gebiet $G$ für die Variable $x$ derart, dass ausserhalb dieses Gebietes der absolute Werth von $f(x)$ immer grösser als $a$ ist, so dass $G$ alle Punkte $x$ enthält, in denen $|f(x)|<a$ ist (aber auch noch andere Punkte). Dies ist nach \S\,31, IV dadurch möglich, dass wir vom Nullpunkt als Mittelpunkt einen die Punkte $\alpha,\beta_1,\beta_2,
\ldots,\beta_{n-1}$ einschliessenden Kreis $(R)$ von hinlänglich grossem Radius $R$ legen und die Punkte $\beta_1,\beta_2,
\ldots,\beta_{n-1}$, in denen $|f(x)|$ ja grösser als $a$ ist, durch kreisförmige Hüllen von so kleinen aber nicht verschwindenden Radien $\rho_1,\rho_2,
\ldots,\rho_{n-1}$ von diesem Kreise ausscheiden, dass im Inneren aller dieser Kreise $(\rho_1),(\rho_2),\ldots,(\rho_{n-1})$ der absolute Werth $|f(x)|$ grösser als $a$ bleibt (\S\,31, V). Dann umschliesst keiner dieser Kreise den Punkt $\alpha$, in dem $|f(x)|=a$ ist. Das von den Kreisen $(R),(\rho_1),\ldots,(\rho_{n-1})$ begrenzte zusammenhängende Flächenstück ist das Gebiet $G$.

Wir bestimmen nun eine Zahlenreihe
\[
 \alpha,\alpha',\alpha'',\alpha''',\ldots
\]
derart, dass die absoluten Werthe von $f(\alpha),f(\alpha'),f(\alpha''),\ldots$, die wir mit
\[
 a,a',a'',a''',\ldots
\]
bezeichnen, immer abnehmen. Die so bestimmten Punkte $\alpha,\alpha',\alpha''$ bleiben alle im Inneren des Gebietes $G$, weil ja in ihnen $|f(x)|<a$ ist.

Wir bedienen uns dazu eines Verfahrens, ganz ähnlich dem zum Beweis des Satzes 5., \S\,38 angewandten, nur dadurch vereinfacht, dass $f'(\alpha)$ schon von Null verschieden ist. Wir setzen, indem wir unter $\delta$ einen noch näher zu bestimmenden, positiven, echten Bruch verstehen, in der Entwickelung
\begin{equation}
 f(\alpha+h)=f(\alpha)+h f'(\alpha)+\frac{h^2}{1\cdot2}f''(\alpha)+\cdots
\tag{6}
\end{equation}
\begin{equation}
 h=-\delta\frac{f(\alpha)}{f'(\alpha)}
\tag{7}
\end{equation}
und erhalten
\begin{align}
 f(\alpha+h)={}&(1-\delta)f(\alpha)
 +\frac{\delta^2 f(\alpha)^2 f''(\alpha)}{1\cdot2\,f'(\alpha)^2}
 -\frac{\delta^3 f(\alpha)^3 f'''(\alpha)}{1\cdot2\cdot3\,f'(\alpha)^3}+\cdots .
\tag{8}
\end{align}

Nun ist nach (2)
\[
 f'(\alpha)=n(\alpha-\beta_1)(\alpha-\beta_2)\cdots(\alpha-\beta_{n-1}),
\]
also, wenn wir die absoluten Werthe von $(\alpha-\beta_1),(\alpha-\beta_2),\ldots,(\alpha-\beta_{n-1})$ durch $r_1,r_2,
\ldots,r_{n-1}$ bezeichnen,
\[
 |f'(\alpha)|=nr_1r_2\cdots r_{n-1}.
\]
Da nun $\alpha$ ausserhalb des um $\beta_i$ beschriebenen Kreises $\rho_i$ liegt, so ist $r_1>\rho_1,r_2>\rho_2,
\ldots$, und folglich
\[
 |f'(\alpha)|>n\rho_1\rho_2\cdots\rho_{n-1},
\]
also $|f'(\alpha)|$ grösser als eine von der Lage von $\alpha$ innerhalb $G$ unabhängige positive Zahl $k$. Daraus ergiebt sich, dass man eine hinlänglich grosse positive Zahl $Q$, die gleichfalls von der Lage von $\alpha$ und von dem echten Bruch $\delta$ unabhängig ist, so wählen kann, dass
\begin{equation}
 \left|\frac{f(\alpha)^2f''(\alpha)}{1\cdot2\,f'(\alpha)^2}
 -\delta\frac{f(\alpha)^3f'''(\alpha)}{1\cdot2\cdot3\,f'(\alpha)^3}+\cdots\right|<Q
\tag{9}
\end{equation}
ist.

Man kann $Q$ erhalten, wenn man auf der linken Seite von (9) in den Nennern $f'(\alpha)$ durch $k$, $\delta$ durch $1$ und sonst alle Glieder durch ihre absoluten Werthe und endlich den absoluten Werth von $\alpha$ durch den grösseren Werth $R$ ersetzt, und die so gewonnene Zahl noch beliebig grösser macht. Dann ergiebt sich aber aus (8)
\begin{equation}
 |f(\alpha+h)|<(1-\delta)|f(\alpha)|+\delta^2Q.
\tag{10}
\end{equation}
Nehmen wir nun $Q$ so gewählt an, dass für jedes $\alpha$ innerhalb $G$
\[
 2Q>|f(\alpha)|=a
\]
ist, was mit der obigen Bestimmung offenbar verträglich ist, so können wir
\begin{equation}
 \delta=\frac{a}{2Q},\qquad
 h=-\frac{a f(\alpha)}{2Q f'(\alpha)}
\tag{11}
\end{equation}
setzen und erhalten, wenn wir $\alpha+h$ mit $\alpha'$ bezeichnen, und wenn wir $|f(\alpha')|$ mit $a'$ bezeichnen, nach (10)
\begin{equation}
 a'<a\left(1-\frac{a}{4Q}\right).
\tag{12}
\end{equation}

Wir leiten nun durch dasselbe Verfahren aus $\alpha',a'$ ein zweites Grössensystem $\alpha'',a''$ ab, dann aus $\alpha'',a''$ ein drittes $\alpha''',a'''$ u. s. f., und erhalten so für die Reihe abnehmender positiver Zahlen $a,a',a'',\ldots$ die Ungleichungen
\begin{equation}
 a'<a\left(1-\frac{a}{4Q}\right),\quad
 a''<a'\left(1-\frac{a'}{4Q}\right),\quad
 a'''<a''\left(1-\frac{a''}{4Q}\right),\ldots .
\tag{13}
\end{equation}
und beweisen nun, dass diese Zahlenreihe sich der Grenze Null nähert. Wäre dies nicht der Fall, so würde sich eine positive Zahl $\omega$ so angeben lassen, dass alle $a^{(\nu)}$ über $\omega$ bleiben; dann wären die Differenzen $a^{(\nu)}/4Q$ alle grösser als der positive echte Bruch $\omega/4Q$, und aus (13) würde folgen
\[
 a'<a\theta,
 \quad a''<a'\theta,
 \quad a'''<a''\theta,
 \quad\ldots,
\]
woraus durch Multiplication
\begin{equation}
 a^{(\nu)}<a\theta^\nu.
\tag{14}
\end{equation}
Darin aber liegt ein Widerspruch, da $\theta^\nu$ mit unendlich wachsendem $\nu$ unendlich klein wird. Das so gewonnene Resultat drücken wir in der Gleichung aus
\begin{equation}
 \lim a^{(\nu)}=0.
\tag{15}
\end{equation}

Wenn wir die Ungleichungen (13) so schreiben
\[
 a-a'>\frac{a^2}{4Q},\qquad
 a'-a''>\frac{a'^2}{4Q},\qquad\ldots,
\]
und sie dann von der $(\mu+1)$ten bis zur $(\nu+1)$ten addiren, so folgt
\[
 (a^{(\mu)})^2+(a^{(\mu+1)})^2+\cdots+(a^{(\nu)})^2
 <4Q\bigl(a^{(\mu)}-a^{(\nu+1)}\bigr),
\]
also, wenn $\mu$ und $\nu$ zwei ins Unendliche wachsende Zahlen bedeuten,
\begin{equation}
 \lim\bigl((a^{(\mu)})^2+(a^{(\mu+1)})^2+\cdots+(a^{(\nu)})^2\bigr)=0.
\tag{16}
\end{equation}

Für die Reihe der Zahlen $\alpha,\alpha',\alpha'',\ldots$ ergiebt sich aus (11) das Bildungsgesetz
\begin{equation}
\begin{aligned}
 \alpha'&=\alpha-\frac{a f(\alpha)}{2Q f'(\alpha)},\\
 \alpha''&=\alpha'-\frac{a' f(\alpha')}{2Q f'(\alpha')},\\
 &\hspace{1.5em}\vdots\\
 \alpha^{(\nu)}&=\alpha^{(\nu-1)}-
 \frac{a^{(\nu-1)} f(\alpha^{(\nu-1)})}{2Q f'(\alpha^{(\nu-1)})}.
\end{aligned}
\tag{17}
\end{equation}
Wenn man eine beliebige Anzahl dieser Gleichungen von der $(\mu+1)$ten bis zur $(\nu+1)$ten addirt, so folgt
\begin{equation}
 \alpha^{(\mu)}-\alpha^{(\nu+1)}=
 \sum_{r=\mu}^{\nu}\frac{a^{(r)}f(\alpha^{(r)})}{2Q f'(\alpha^{(r)})}.
\tag{18}
\end{equation}
Daraus ergiebt sich, da die sämmtlichen $|f'(\alpha^{(r)})|\ge k$ sind, für den absoluten Werth dieser Differenz
\begin{equation}
 |\alpha^{(\mu)}-\alpha^{(\nu+1)}|<
 \frac{(a^{(\mu)})^2+(a^{(\mu+1)})^2+\cdots+(a^{(\nu)})^2}{2Qk},
\tag{19}
\end{equation}
also mit Hülfe von (16)
\begin{equation}
 \lim |\alpha^{(\mu)}-\alpha^{(\nu+1)}|=0.
\tag{20}
\end{equation}

Damit ist ausgedrückt (nach dem in der Einleitung erklärten Begriff der Zahlenreihe, S. 15), dass sich die Zahlen $\alpha,\alpha',\alpha'',\ldots$ einer bestimmten Grenze nähern, die wir mit $\xi$ bezeichnen wollen. Aus der Stetigkeit der Function $f(x)$ folgt aber dann leicht, dass $f(\xi)=0$ sein muss; denn wäre $|f(\xi)|>0$, so könnten, da die $\alpha^{(\nu)}$ dem Werthe $\xi$ beliebig nahe kommen, die absoluten Werthe $a^{(\nu)}$ von $f(\alpha^{(\nu)})$ nicht unter jeden positiven Werth herunter sinken, was wir doch von ihnen nachgewiesen haben.

\end{document}
